\documentclass[10pt]{extarticle}
\usepackage[a4paper,margin=1.5cm]{geometry}
\usepackage[T1]{fontenc}
\usepackage{amsmath,amssymb}
\usepackage{graphicx}
\usepackage[export]{adjustbox}
\usepackage{multicol}
\usepackage{enumitem}
\usepackage{booktabs,array,tabularx}
\setlength{\columnsep}{0.7cm}
\setlength{\parindent}{0pt}
\setlength{\parskip}{3pt}
\pagestyle{empty}
\makeatletter
\newcommand{\qfigdef}[4]{\@namedef{qf@file@#1}{#2}\@namedef{qf@w@#1}{#3}\@namedef{qf@wide@#1}{#4}}
\newcommand{\qf@check}[1]{\@ifundefined{qf@file@#1}{\PackageError{score2}{Unknown figure id #1}{}}{}}
\newcommand{\qf@img}[4][]{\includegraphics[width=#2\dimexpr\@nameuse{qf@w@#4}pt\relax,max width=#3,max height=0.45\textheight,#1]{\@nameuse{qf@file@#4}}}
\newif\ifqf@full
\newif\ifqf@boxed
\newsavebox{\qf@box}
\newcommand{\qf@open}{\global\qf@fulltrue\par\global\qf@boxedtrue
  \setbox\qf@box\vbox\bgroup\hsize\textwidth\linewidth\textwidth}
\newcommand{\qf@emit}[1]{\par\egroup\global\qf@boxedfalse
  \ifdim\dimexpr\ht\qf@box+\dp\qf@box\relax<#1\textheight\noindent\usebox{\qf@box}\par\else\unvbox\qf@box\fi}
\newcommand{\qfullsplit}{\ifqf@boxed\qf@emit{0.85}\fi}
\newcommand{\qf@close}{\ifqf@boxed\qf@emit{0.45}\fi\global\qf@fullfalse}
\newcommand{\qfullbegin}{\end{multicols}\qf@open}
\newcommand{\qfullend}{\qf@close\begin{multicols}{2}}
\newcommand{\qfullstart}{\qf@open}
\newcommand{\qfulldone}{\qf@close}
\newcommand{\qfig}[2][1]{\qf@check{#2}%
  \ifqf@full\qf@blk{#1}{#2}\else
  \ifnum\@nameuse{qf@wide@#2}=1 \expandafter\@firstoftwo\else\expandafter\@secondoftwo\fi
  {\end{multicols}\begin{center}\qf@img{#1}{\textwidth}{#2}\end{center}\begin{multicols}{2}}%
  {\qf@blk{#1}{#2}}\fi}
\newcommand{\qf@blk}[2]{\par\smallskip{\centering\qf@img{#1}{\linewidth}{#2}\par}\smallskip}
\newcommand{\qfigi}[2][1]{\qf@check{#2}\qf@img[valign=c]{#1}{\linewidth}{#2}}
\makeatother
\newcommand{\papersec}[1]{\par\bigskip{\centering\Large\bfseries PAPER-#1\par}\medskip}
\newcommand{\partsec}[2]{\par\medskip\noindent{\large\bfseries PART-#1: #2}\par\smallskip}
\newcommand{\sect}[1]{\par\smallskip\noindent\textbf{#1}\par}
\newcommand{\qn}[1]{\par\smallskip\noindent\textbf{#1.}~}
\newcommand{\paraq}[2]{\par\smallskip\noindent\textbf{Paragraph for Questions #1 and #2}\par}
\newenvironment{qoptions}{\begin{itemize}[leftmargin=1.8em,itemsep=1pt,topsep=2pt,label={}]}{\end{itemize}}
\newcommand{\opt}[2]{\item[(#1)] #2}
% figure registry (generated by score2_tools.py finalize)
\qfigdef{p1_phy_s1_q01_1}{Score-I - Test 06 (28 Jan 2022)-img/p1_phy_s1_q01_1.png}{131.8}{0}
\qfigdef{p1_phy_s1_q02_1}{Score-I - Test 06 (28 Jan 2022)-img/p1_phy_s1_q02_1.png}{282.6}{1}
\qfigdef{p1_phy_s1_q05_1}{Score-I - Test 06 (28 Jan 2022)-img/p1_phy_s1_q05_1.png}{138.6}{0}
\qfigdef{p1_phy_s1_q06_1}{Score-I - Test 06 (28 Jan 2022)-img/p1_phy_s1_q06_1.png}{150.1}{0}
\qfigdef{p1_phy_s2_q07_1}{Score-I - Test 06 (28 Jan 2022)-img/p1_phy_s2_q07_1.png}{206.1}{0}
\qfigdef{p1_phy_s2_q09_1}{Score-I - Test 06 (28 Jan 2022)-img/p1_phy_s2_q09_1.png}{104.3}{0}
\qfigdef{p1_phy_s2_q11_1}{Score-I - Test 06 (28 Jan 2022)-img/p1_phy_s2_q11_1.png}{123.0}{0}
\qfigdef{p1_phy_s3_q02_1}{Score-I - Test 06 (28 Jan 2022)-img/p1_phy_s3_q02_1.png}{180.0}{0}
\qfigdef{p1_phy_s3_q03_1}{Score-I - Test 06 (28 Jan 2022)-img/p1_phy_s3_q03_1.png}{149.6}{0}
\qfigdef{p1_phy_s3_q04_1}{Score-I - Test 06 (28 Jan 2022)-img/p1_phy_s3_q04_1.png}{149.6}{0}
\qfigdef{p1_phy_s3_q05_1}{Score-I - Test 06 (28 Jan 2022)-img/p1_phy_s3_q05_1.png}{205.7}{0}
\qfigdef{p1_phy_s3_q06_1}{Score-I - Test 06 (28 Jan 2022)-img/p1_phy_s3_q06_1.png}{316.2}{1}
\qfigdef{p1_chem_s1_q05_1}{Score-I - Test 06 (28 Jan 2022)-img/p1_chem_s1_q05_1.png}{283.1}{1}
\qfigdef{p1_chem_s1_q06_A}{Score-I - Test 06 (28 Jan 2022)-img/p1_chem_s1_q06_A.png}{124.2}{0}
\qfigdef{p1_chem_s1_q06_B}{Score-I - Test 06 (28 Jan 2022)-img/p1_chem_s1_q06_B.png}{95.4}{0}
\qfigdef{p1_chem_s1_q06_C}{Score-I - Test 06 (28 Jan 2022)-img/p1_chem_s1_q06_C.png}{98.9}{0}
\qfigdef{p1_chem_s1_q06_D}{Score-I - Test 06 (28 Jan 2022)-img/p1_chem_s1_q06_D.png}{87.5}{0}
\qfigdef{p1_chem_s2_q08_1}{Score-I - Test 06 (28 Jan 2022)-img/p1_chem_s2_q08_1.png}{116.1}{0}
\qfigdef{p1_chem_s3_q05_i}{Score-I - Test 06 (28 Jan 2022)-img/p1_chem_s3_q05_i.png}{76.5}{0}
\qfigdef{p1_chem_s3_q05_ii}{Score-I - Test 06 (28 Jan 2022)-img/p1_chem_s3_q05_ii.png}{124.4}{0}
\qfigdef{p1_chem_s3_q05_iii}{Score-I - Test 06 (28 Jan 2022)-img/p1_chem_s3_q05_iii.png}{52.4}{0}
\qfigdef{p1_chem_s3_q05_iv}{Score-I - Test 06 (28 Jan 2022)-img/p1_chem_s3_q05_iv.png}{39.9}{0}
\qfigdef{p1_chem_s3_q05_v}{Score-I - Test 06 (28 Jan 2022)-img/p1_chem_s3_q05_v.png}{81.3}{0}
\qfigdef{p1_chem_s3_q05_vi}{Score-I - Test 06 (28 Jan 2022)-img/p1_chem_s3_q05_vi.png}{118.3}{0}
\qfigdef{p1_chem_s3_q06_1}{Score-I - Test 06 (28 Jan 2022)-img/p1_chem_s3_q06_1.png}{126.3}{0}
\qfigdef{p1_math_s3_q02_1}{Score-I - Test 06 (28 Jan 2022)-img/p1_math_s3_q02_1.png}{80.8}{0}

\begin{document}
\begin{center}{\Large\bfseries JEE (Main+Advanced) Score-I --- Test 06 (28 Jan 2022)}\end{center}
\begin{multicols}{2}
\raggedcolumns
\partsec{1}{PHYSICS}
\sect{Section-I(i)}
\qn{1} A conducting medium is shaped in a form of a quarter of an annulus of radii b and a (b $>$ a) and thickness t. Specific resistance of the medium is $\rho$. Find the ratio of resistance between faces LMNO and PQRS to resistance between faces LMRS and OPQN.
\qfig{p1_phy_s1_q01_1}
\begin{qoptions}
\opt{A}{$\dfrac{2}{\pi^2}\left[\ln\left(\dfrac{b}{a}\right)\right]^2$}
\opt{B}{$\dfrac{4(b-a)}{\pi^2 t}\left[\ln\left(\dfrac{b}{a}\right)\right]^2$}
\opt{C}{$\dfrac{4}{\pi^2}\left[\ln\left(\dfrac{b}{a}\right)\right]^2$}
\opt{D}{$\dfrac{4t}{\pi^2(b-a)}\left[\ln\left(\dfrac{b}{a}\right)\right]^2$}
\end{qoptions}
\qfullbegin
\qn{2} Maximum rate of chemical energy consumed by the battery in watt is $P_0$. Find $\dfrac{P_0}{10}$. (EMF of the battery is E = 10V and its internal resistance is r = 0.5 $\Omega$). R is choosen such that rate of chemical energy consumed by battery is maximum.
\qfig{p1_phy_s1_q02_1}
\qfullsplit
\begin{qoptions}
\opt{A}{10}
\opt{B}{15}
\opt{C}{12}
\opt{D}{20}
\end{qoptions}
\qfullend
\qn{3} The main scale of vernier callipers reads in millimetre and its vernier is divided into 8 divisions, which coincide with 5 divisions of main scale. When two jaws of instrument touch each other the zero of the vernier coincide with the zero of main scale. A rod is tightly placed along its length between both jaws, it is observed that the zero vernier scale lies just left to 36$^\text{th}$ division of main scale and fourth division of vernier scale coincide with the main scale. Then the measured value is --
\begin{qoptions}
\opt{A}{3.66 cm}
\opt{B}{3.55 cm}
\opt{C}{3.65 cm}
\opt{D}{3.56 cm}
\end{qoptions}
\qn{4} Complete the following calculations, reporting your answer to the correct number of significant figures :
\[\dfrac{10{,}569-958.82-1.8\times10^{-1}}{0.9950+8.60+0.015}\]
\begin{qoptions}
\opt{A}{1000}
\opt{B}{$1\times10^3$}
\opt{C}{$1.00\times10^3$}
\opt{D}{All of these}
\end{qoptions}
\qn{5} Five identical capacitor (initially uncharged) are connected in circuit as shown in figure. Capacitance of each capacitor is C. In steady state charge on capacitor `1' is
\qfig{p1_phy_s1_q05_1}
\begin{qoptions}
\opt{A}{$\dfrac{14C\,\varepsilon}{45}$}
\opt{B}{$\dfrac{4C\,\varepsilon}{9}$}
\opt{C}{$\dfrac{2C\,\varepsilon}{3}$}
\opt{D}{Zero}
\end{qoptions}
\qn{6} Consider a cuboidal vessel $(2R\times2R\times R)$ with a hemispherical cavity of radius R, kept at a horizontal smooth surface as shown in the figure. The vessel has very small hole of cross-sectional area ``a''. Now water of density $\rho$ is poured into space developed due horizontal surface and vessel through hole very slowly. When the height of water level is h, the vessel lifts off the surface and liquid leaks through space generated. Find the value of h. (R = 2m, $\rho=10^3$ kg/m$^3$, mass of vessel m = 9$\pi$ kg); (Atmospheric pressure = $10^5$ N/m$^2$).
\qfig{p1_phy_s1_q06_1}
\begin{qoptions}
\opt{A}{20 cm}
\opt{B}{15 cm}
\opt{C}{25 cm}
\opt{D}{30 cm}
\end{qoptions}
\sect{Section-I(ii)}
\qn{7} A solid cone of density d has $\dfrac13$rd of its length in a fluid of density d$_1$, other $\dfrac13$rd in air and the remaining $\dfrac13$rd in a fluid of density d$_2$, as shown in figure. Given that d$_1=\dfrac35$d ; the cone is in equilibrium in position shown with the contact force from the upper container's bottom being negligibly small, atmospheric pressure due to air is P$_0$ and force by atmospheric pressure on the middle part in air is F. Pick correct option(s)
\qfig{p1_phy_s2_q07_1}
\begin{qoptions}
\opt{A}{d$_2$ = 10 d}
\opt{B}{F = P$_0\times2\pi$h$^2$}
\opt{C}{F = P$_0\times3\pi$h$^2$}
\opt{D}{d$_2$ = 30d}
\end{qoptions}
\qn{8} A water drop and a mercury drop is taken inside identical horizontal conical glass pipes. Then
\begin{qoptions}
\opt{A}{The water drop tends to move towards the narrow end.}
\opt{B}{The water drop tends to move towards the wide end.}
\opt{C}{The mercury drop tends to move towards the wide end.}
\opt{D}{The mercury drop tends to move towards the narrow end.}
\end{qoptions}
\qn{9} Consider the circuit shown. The resistance connected between the junction A and B is 60 $\Omega$ including the resistance of the galvanometer. The switches have no resistance when shorted and infinite resistance when opened. All the switches are initially open and they are closed as given in the options. Choose correct option(s) :
\qfig{p1_phy_s2_q09_1}
\begin{qoptions}
\opt{A}{If only switch S$_1$ is closed current through the battery is 20 Amp}
\opt{B}{If only switch S$_2$ is closed then current in galvanometer from B to A}
\opt{C}{If only switch S$_3$ is closed then current in galvanometer from B to A}
\opt{D}{If only switch S$_4$ is closed then current in galvanometer from A to B}
\end{qoptions}
\qn{10} We have five identical galvanometer of resistance R and five identical resistance R.

Situation $\to$ (1) all galvanometer and resistance are connected in series.

Situation $\to$ (2) all galvanometer connected in parallel, all resistance connected in parallel and combination of galvanometers and combination of resistances connected in series.

Situation $\to$ (3) all galvanometer connected in parallel, all resistance connected in series and combination of galvanometers and resistances are connected in series.

Situation $\to$ (4) all galvanometer connected in series, all resistance connected in parallel and combination of galvanometers and resistance connected in series.

The complete combination in all situation acts as voltmeter.
\begin{qoptions}
\opt{A}{Range of voltmeter in situation (1) is 5 times range in situation (2) and $>$ range in situation (4)}
\opt{B}{Range of voltmeter in situation (1) is 5 times range in situation (2) and $<$ range in situation (4)}
\opt{C}{The value of range of voltmeter is maximum in the situation (3). (In above all given situation)}
\opt{D}{Range of voltmeter in situation (3) is $>$ range in situation (4)}
\end{qoptions}
\qn{11} Two identical capacitor A and B of capacitance `C' each connected in the circuit as shown in the figure with ideal battery. Capacitor B having charge Q$_0$ = C$\varepsilon$ initially. At t = 0 both S$_1$ and S$_2$ closed simultaneously.
\qfig{p1_phy_s2_q11_1}
\begin{qoptions}
\opt{A}{Current in R as function of time is $\dfrac{\varepsilon}{R}e^{-t/2RC}$}
\opt{B}{Charge on capacitor A as function of time is C$\varepsilon\left(1-e^{-t/2RC}\right)$}
\opt{C}{Charge on capacitor B in steady state is zero}
\opt{D}{Total work done by the battery till steady state is reached is 2C$\varepsilon^2$}
\end{qoptions}
\qn{12} The universal gravitational constant is measured experimentally using Cavendish's experiment. This experiment is a very precise experiment which gives the value of G as $(6.6742\pm0.0010)\times10^{-11}$ N.m$^2$/kg$^2$. Since the universal gravitational constant has some error in it, using it if we find escape speed and acceleration due to gravity on Earth's surface, our values will have error in it. (Assume mass of Earth = $6.0000\times10^{24}$ kg \& radius of Earth = $6.4000\times10^6$ m.) What will be the maximum percentage error in calculating the escape speed on Earth ? (The measurement of mass of Earth has an error of 0.01 \% and radius of Earth has an error of 0.02\% ):-
\begin{qoptions}
\opt{A}{0.045\%}
\opt{B}{0.0225 \%}
\opt{C}{0.015 \%}
\opt{D}{0.0025\%}
\end{qoptions}
\sect{Section-II}
\qn{1} The pitch of screw gauge is 1 mm and its circular scale is divided into 50 divisions. When nothing is put between the studs 44$^\text{th}$ division of circular scale coincides with the reference line and circular scale lies between 3$^\text{rd}$ and 4$^\text{th}$ division of main scale. When a glass plate is placed between the studs, the circular scale lies between 13$^\text{th}$ and 14$^\text{th}$ division of main scale and circular scale reads 26 divisions. Find the thickness of the plate (in mm).
\qn{2} Two balls of radii r and $\dfrac{r}{2}$ are released inside a deep water tank. Their initial accelerations are found to be $\dfrac{g}{2}$ and $\dfrac{g}{4}$ downward respectively. If the velocity of smaller ball relative to the larger ball, a long time after the two balls are released is $\dfrac{x\,r^2\rho g}{54\,\eta}$. Coefficient of viscosity is given to be $\eta$. Find x. ($\rho$ = density of water)
\qfig{p1_phy_s3_q02_1}
\qn{3} Two vary large open tanks A and F contain the same liquid as shown in the figure. A horizontal pipe BCD having a constriction at C and open to the air at D, leads out of the bottom of tank A, and a vertical pipe E opens into the constriction at C and dips into the liquid in tank F. Assume streamline flow and no viscosity. If the cross sectional area at C is one half the area at D and if D is a distance 3 cm below the level of the liquid in A, upto what height h$_2$ (cm) will the liquid rise in pipe E? (Neglect surface tension)
\qfig{p1_phy_s3_q03_1}
\qn{4} In the given circuit diagram, if ideal ammeter connected between point A and B its reading is 5A, if ammeter of resistance 3$\Omega$ is connected between A and B its reading is 3A. What will be reading of ideal voltmeter if it is connected between A and B (in volt)?
\qfig{p1_phy_s3_q04_1}
\qn{5} In the given circuit diagram initially all capacitor are uncharged. Answer the following question in steady state. Heat developed during charging is nC$\varepsilon^2$. Find $\dfrac1n$.
\qfig{p1_phy_s3_q05_1}
\qfullbegin
\qn{6} In the circuit shown (figure 1), all the switches are open initially and all capacitors are uncharged. Now first only S$_1$ is closed and then opened when steady state is reached. Then S$_2$ is closed and again opened when steady state is reached, after this S$_3$ is closed and opened when again steady state is reached, this process repeated countineously. After closing sixth switch, the fourth capacitor (8C) is removed and connected (say at t = 0) to an uncharged capacitor of original capacitance C which now has a dielectric (dielectric constant 2) filling half its volume (see figure 2), through a resistor R. If final charge on capacitor B when steady state reached in figure 2 is $\dfrac{8}{n}$CV. The value of n is
\qfig{p1_phy_s3_q06_1}
\qfullsplit
\raggedcolumns
\qfullend
\partsec{2}{CHEMISTRY}
\sect{Section-I(i)}
\qn{1} Select incorrect statement wrt H (protium), D and T
\begin{qoptions}
\opt{A}{Both D \& T can be used for radioactive--labelling in organic compound.}
\opt{B}{Tritium decays by $\beta$--emission.}
\opt{C}{All three can't be used for radioactive--labelling in organic compound.}
\opt{D}{Half-life of T is approximate 12.3 years.}
\end{qoptions}
\qn{2} Which of the following pair of reagents can be used for producing hydrogen gas?

(I) Zn + dil. H$_2$SO$_4$ \quad (II) Zn + NaOH \quad (III) Al + HCl (g) \quad (IV) Al + NaOH
\begin{qoptions}
\opt{A}{I, II and III only}
\opt{B}{II, III and IV only}
\opt{C}{I, III and IV only}
\opt{D}{I, II, III and IV}
\end{qoptions}
\qn{3} For a gaseous chemical reaction the value of K$_\mathrm{P}$ = $\dfrac43$atm$^{-1}$ at 227$^\circ$C in a closed container. The correct statement for the reaction at same temperature is/are : (Given R = $\dfrac{1}{12}\dfrac{\text{litre}-\text{atm}}{\text{k}-\text{mol}}$)
\begin{qoptions}
\opt{A}{K$_\mathrm{C}$ = $\dfrac{20}{3}$ mole litre$^{-1}$}
\opt{B}{If volume of container is increased equilibrium shifts in backward direction.}
\opt{C}{If partial pressure of each reactant is doubled K$_\mathrm{P}$ becomes 4 times.}
\opt{D}{Can't measure K$_\mathrm{C}$ for the reaction as data are insufficient.}
\end{qoptions}
\qn{4} What should the solubility of AgCN in acidic buffer solution of pH = 5.

Ka$_{(\mathrm{HCN})}$ = $5\times10^{-10}$ ; K$_{\mathrm{sp}\,(\mathrm{AgCN})}$ = $6\times10^{-17}$
\begin{qoptions}
\opt{A}{$7.7\times10^{-9}$ M}
\opt{B}{$1.1\times10^{-6}$ M}
\opt{C}{$2.1\times10^{-7}$ M}
\opt{D}{$3.2\times10^{-8}$ M}
\end{qoptions}
\qfullbegin
\qn{5} The correct stability order of the following species is.
\qfig{p1_chem_s1_q05_1}
\qfullsplit
\begin{qoptions}
\opt{A}{c $<$ a $<$ b}
\opt{B}{c = b $<$ a}
\opt{C}{c $<$ a = b}
\opt{D}{a = b = c}
\end{qoptions}
\qfullend
\qn{6} Which of the following is incorrect for Enol \% ?
\begin{qoptions}
\opt{A}{\qfigi{p1_chem_s1_q06_A}}
\opt{B}{\qfigi{p1_chem_s1_q06_B}}
\opt{C}{\qfigi{p1_chem_s1_q06_C}}
\opt{D}{\qfigi{p1_chem_s1_q06_D}}
\end{qoptions}
\sect{Section-I(ii)}
\qn{7} Which of the following incorrect for stereoisomers 2-chloro-1-fluoro propene.
\begin{qoptions}
\opt{A}{E--isomer is having more water solubility then Z--isomer.}
\opt{B}{Z--isomer is having more water solubility then E--isomer}
\opt{C}{E \& Z both isomers are equally water soluble}
\opt{D}{E--isomer is having intramolecular H--bonding.}
\end{qoptions}
\qn{8} Correct R, S configuration for following compound is:
\qfig{p1_chem_s2_q08_1}
\begin{qoptions}
\opt{A}{R at A, R at B}
\opt{B}{R at A, S at B}
\opt{C}{S at A, R at B}
\opt{D}{S at A, S at B}
\end{qoptions}
\qn{9} Identify the options which is/are correct w.r.t properties of elements.
\begin{qoptions}
\opt{A}{IE of Li$^-$ will be more than Na$^-$}
\opt{B}{Inspite of lower charge, the ionic potential (charge/size ratio) of Li$^+$ is more than that of Mg$^{2+}$}
\opt{C}{Cs is the best reducing agent, therefore among alkali metal, reducing power is influenced by hydration energy.}
\opt{D}{Correct order of EN of 13$^\text{th}$ group elements is B $>$ T$\ell$ $>$ In $>$ Ga $>$ Al.}
\end{qoptions}
\qn{10} Match the following.

In List--I the comparision of IE values and EA values of X and Y are given and in List--II. The elements X and Y are listed. Match them.
\begin{center}\small\setlength{\tabcolsep}{3pt}
\begin{tabular}{|c|p{3.6cm}|c|c|c|}
\hline
 & List I & & \multicolumn{2}{c|}{List II}\\ \hline
 & & & X & Y\\ \hline
(P) & IE$_1$ of X is more than that of Y & (1) & Be & B\\ \hline
(Q) & IE$_2$ of X is more than that of Y & (2) & B & C\\ \hline
(R) & IE$_3$ of X is more than that of Y & (3) & C & N\\ \hline
(S) & EA of X is less than that of Y & (4) & N & O\\ \hline
\end{tabular}
\end{center}
\begin{qoptions}
\opt{A}{P $\to$ 1,4; Q $\to$ 2; R $\to$ 1,3; S $\to$ 1,2,4}
\opt{B}{P $\to$ 1,2,4; Q $\to$ 2; R $\to$ 1,3; S $\to$ 1,2,4}
\opt{C}{P $\to$ 1,3; Q $\to$ 1,2; R $\to$ 1,3; S $\to$ 1,2,4}
\opt{D}{P $\to$ 2,3; Q $\to$ 1,2; R $\to$ 1,3; S $\to$ 1,2,4}
\end{qoptions}
\qn{11} Following two equilibria are simultaneously established in a container

$\mathrm{PCl_5(g)\rightleftharpoons PCl_3(g)+Cl_2}$ ; $\mathrm{CO(g)+Cl_2(g)\rightleftharpoons COCl_2(g)}$

If some Ni(s) is introduced in the container establishing another equilibrium.

Which among the following statements is/are definitely true about this system ?
\begin{qoptions}
\opt{A}{COCl$_2$ concentration decreases as the new equilibrium is established.}
\opt{B}{As Ni is a solid it will not have any impact on the existing two equilibrium reactions.}
\opt{C}{PCl$_3$ concentration decreases to re--attain the equilibrium.}
\opt{D}{PCl$_5$ concentration increases to attain the new equilibrium.}
\end{qoptions}
\qn{12} Solubility product of AgCl (s) in water at 25$^\circ$C is $12\times10^{-10}$ M$^2$ then select the correct options.
\begin{qoptions}
\opt{A}{Solubility of AgCl in presence of $10^{-5}$ M AgNO$_3$ is $3\times10^{-5}$ moles/L.}
\opt{B}{Solubility of AgCl in presence of excess of KCl will increase due to formation of complex AgCl$_2^-$(aq.)}
\opt{C}{Solubility of AgCl in presence of 2M NH$_3$ is greater than that in pure water.}
\opt{D}{Solubility of AgCl remains unaffected on adding any amount of NaBr because there is no common ion effect.}
\end{qoptions}
\sect{Section-II}
\qn{1} 0.005 mol of H$_2$A$^-$ and 0.105 mol A$^{3-}$ are mixed in 1 litre solution. Find out the pH of final solution. (H$_3$A Ka$_1=10^{-5}$, Ka$_2=10^{-9}$ Ka$_3=10^{-12}$)
\qn{2} $\mathrm{HCN\rightleftharpoons H^++CN^-}$ \hfill ; K = $10^{-9}$

$\mathrm{NH_4OH\rightleftharpoons NH_4+OH^-}$ \hfill ; K = $10^{-5}\times2$

$\mathrm{HCN+KOH\to KCN+H_2O}$ ; K$_1$

$\mathrm{NH_4OH+HCN\to NH_4CN+H_2O}$ ; K$_2$

$\mathrm{NH_4NO_3+H_2O\to NH_4OH+HNO_3}$ + K$_3$

$\dfrac{\mathrm{K_1\times K_3}}{\mathrm{K_2}}=\mathrm{x}\times10^{-7}$ \ ; `x' is :
\qn{3} In how many cases electron gain enthalpy is negative.

(i) \ $\mathrm{H(g)+e\to H^-(g)}$

(ii) $\mathrm{Li(g)+e\to Li^-(g)}$

(iii) $\mathrm{Na(g)+e\to Na^-(g)}$

(iv) $\mathrm{K(g)+e\to K^-(g)}$

(v) $\mathrm{P(g)+e\to P^-(g)}$

(vi) $\mathrm{O(g)+e\to O^-(g)}$

(vii) $\mathrm{O^-(g)+e\to O^{2-}(g)}$

(viii) $\mathrm{N(g)+e\to N^-(g)}$

(ix) $\mathrm{Ar(g)+e\to Ar^-(g)}$
\qn{4} How many incorrect statements ?

(a) CO which is a major pollutant resulting from the combustion of fuels in automobiles plays a major role in photochemical smog.

(b) Classical smog has an oxidising character while the photochemical smog is reducing in character.

(c) Photochemical smog occurs in day time whereas the classical smog occurs in the morning hours.

(d) During formation of smog the level of ozone in the atmosphere goes down.

(e) Classical smog is good for health but not photochemical smog.

(f) Chlorofluorocarbons are responsible for ozone layer depletion.

(g) Green house effect is responsible for global warming.

(h) Ozone layer does not permit infrared radiation from the sum to reach the earth.

(i) Acid rain is mostly because of oxides of nitrogen and sulphur.
\qn{5} How many of the following compounds are chiral ?

(i) \qfigi{p1_chem_s3_q05_i}

(ii) \qfigi{p1_chem_s3_q05_ii}

(iii) \qfigi{p1_chem_s3_q05_iii}

(iv) \qfigi{p1_chem_s3_q05_iv}

(v) \qfigi{p1_chem_s3_q05_v}

(vi) \qfigi{p1_chem_s3_q05_vi}
\qn{6} \qfigi{p1_chem_s3_q06_1}

For above compound's

P $\to$ total number of stereoisomers.

Q $\to$ total number of enantiomeric pairs

R $\to$ total number of meso compounds

Then give your answer in 3 digit's as PQR.
\raggedcolumns
\partsec{3}{MATHEMATICS}
\sect{Section-I(i)}
\qn{1} The interval to which b must belong so that the function $f(x)=\left(1-\dfrac{\sqrt{21-4b-b^2}}{b+1}\right)x^3+5x+\sqrt6$ is increasing at every point of its domain is
\begin{qoptions}
\opt{A}{$[-10,-1)$}
\opt{B}{$[-8,-6]$}
\opt{C}{$[2,3]$}
\opt{D}{$[3,8]$}
\end{qoptions}
\qn{2} If $\displaystyle\int u\dfrac{d^2v}{dx^2}dx=u\dfrac{dv}{dx}-v\dfrac{du}{dx}+w$ then w is equal to :
\begin{qoptions}
\opt{A}{$\displaystyle\int v\dfrac{d^2u}{dx^2}dx$}
\opt{B}{$\displaystyle\int u\dfrac{d^2v}{dx^2}dx$}
\opt{C}{$\displaystyle\int v\left(\dfrac{du}{dx}\right)^2dx$}
\opt{D}{$\displaystyle\int u\left(\dfrac{dv}{dx}\right)^2dx$}
\end{qoptions}
\qn{3} Which of the following equation has real solution(s) ?
\begin{qoptions}
\opt{A}{$\cos(\sin x)=\dfrac12$}
\opt{B}{$(x^2+1)^{10}=2x-x^2-2$}
\opt{C}{$(3+\cos x)^2=4-2\sin^2x$}
\opt{D}{$\sin x+\cos x=\dfrac{\pi}{2}$}
\end{qoptions}
\qn{4} There is a point P inside an equilateral triangle ABC of side d whose distances from the vertices A,B \& C are 3, 4, 5 respectively. Rotate the triangle through $60^\circ$ about C. Let A goes to A$'$ and P to P$'$. The area of triangle PAP$'$ is
\begin{qoptions}
\opt{A}{8}
\opt{B}{12}
\opt{C}{6}
\opt{D}{None of these}
\end{qoptions}
\qn{5} Tangent is drawn to ellipse $\dfrac{x^2}{27}+y^2=1$ at $\left(3\sqrt3\cos\theta,\sin\theta\right)$ $\left(\text{where }\theta\in\left(0,\dfrac{\pi}{2}\right)\right)$. Then the value of $\theta$ such that sum of intercepts on axes made by this tangent is least, is
\begin{qoptions}
\opt{A}{$\dfrac{\pi}{3}$}
\opt{B}{$\dfrac{\pi}{6}$}
\opt{C}{$\dfrac{\pi}{8}$}
\opt{D}{$\dfrac{\pi}{4}$}
\end{qoptions}
\qn{6} If f : R $\to$ R is a function such that $f(x)=x^3+x^2f'(1)+xf''(2)+f'''(3)$ for $x\in\mathbb{R}$ then the value of $f(2)$ is
\begin{qoptions}
\opt{A}{5}
\opt{B}{10}
\opt{C}{6}
\opt{D}{$-2$}
\end{qoptions}
\sect{Section-I(ii)}
\qn{7} If a, b and c be in A.P. Then the correct results from the given statements are :
\begin{qoptions}
\opt{A}{$\cot\dfrac{A}{2}$, $\cot\dfrac{B}{2}$ and $\cot\dfrac{C}{2}$ are in A.P.}
\opt{B}{$\cos A\cot\dfrac{A}{2}$, $\cos B\cot\dfrac{B}{2}$ and $\cos C\cot\dfrac{C}{2}$ are in A.P.}
\opt{C}{$a\cos^2\dfrac{C}{2}+c\cos^2\dfrac{A}{2}=\dfrac{3b}{2}$}
\opt{D}{$\cot\dfrac{A}{2}\cot\dfrac{C}{2}=2$}
\end{qoptions}
\qn{8} If $y=mx+5$ is a tangent to the curve $x^3y^3=ax^3+by^3$ at P(1, 2), then
\begin{qoptions}
\opt{A}{$a+b=\dfrac{19}{5}$}
\opt{B}{$a>b$}
\opt{C}{$a<b$}
\opt{D}{$a+b=\dfrac{5}{18}$}
\end{qoptions}
\qn{9} $\dfrac{\alpha^2}{1-\tan^2x}=\dfrac{\sin^2x+\alpha^2-2}{\cos2x}$ has a solution if
\begin{qoptions}
\opt{A}{$\alpha\in(-\infty,-1)\sim\{-\sqrt3\}$}
\opt{B}{$\alpha\in(1,\infty)\sim\{\sqrt3\}$}
\opt{C}{$\alpha\in(-1,1)$}
\opt{D}{$\alpha$ is any real number}
\end{qoptions}
\qn{10} Let $f(x)=\displaystyle\int_0^x e^{-t^2}(t-5)(t^2-7t+12)\,dt$ for all $x\in(0,\infty)$, then
\begin{qoptions}
\opt{A}{f has a local maximum at x = 4 and a local minimum at x = 3.}
\opt{B}{f is decreasing on $(3,4)\cup(5,\infty)$ and increasing on $(0,3)\cup(4,5)$.}
\opt{C}{There exists atleast two $c\in(0,\infty)$ such that $f''(c)=0$}
\opt{D}{There exists some $c\in(0,\infty)$ such that $f'''(c)=0$}
\end{qoptions}
\qn{11} Let $f(x)=e^{ax}\sin(bx+c)$ and $f''(x)=r^2e^{ax}\sin(bx+\theta)$, then (a,b,c,r and $\theta$ are real positive arbitrary constants and b $<$ a)
\begin{qoptions}
\opt{A}{$r=a^2+b^2$}
\opt{B}{$r=\sqrt{a^2+b^2}$}
\opt{C}{$\theta=c+2\tan^{-1}\left(\dfrac{b}{a}\right)$}
\opt{D}{$\theta=2\mathrm{atan}^{-1}\left(\dfrac{b}{a}\right)$}
\end{qoptions}
\qn{12} Let f be a continuous function and $f''(x)<0\ \forall\ x\in\mathbb{R}$. Let $f(0)=1$ and $f(1)=0$. Which of the following statements are necessarily true ?
\begin{qoptions}
\opt{A}{There exists a number k in [0, 1] such that $f(k)\geq f(x)$ for all x in [0, 1]}
\opt{B}{$f'(1)<-1$}
\opt{C}{There exists a number k in [0, 1] such that $f(k)=\dfrac12$}
\opt{D}{$f'(0)>-1$}
\end{qoptions}
\sect{Section-II}
\qn{1} The number of solutions that the equation $\sin(\cos(\sin x))=\cos(\sin(\cos x))$ has in $\left[0,\dfrac{\pi}{2}\right]$ is :
\qn{2} In the adjoining figure, AC = BC, $\angle$ACB = $96^\circ$ and D is a point inside $\triangle$ABC such that $\angle$DAB = $18^\circ$, $\angle$DBA = $30^\circ$. Then the value of $\left[\dfrac{\angle\mathrm{ADC}}{17}\right]$ is
\qfig{p1_math_s3_q02_1}
(where [.] denotes greater integer function)
\qn{3} Let $y=f(x)$ be a one-one function satisfying the equation $\dfrac{d^2y}{dx^2}+\dfrac{d^2x}{dy^2}=0$ and $f'(0)=1$, $f(1)=2$, then $f(0)$ is :
\qn{4} If product of all possible values of k for which angle of intersection of curves $y=e^{3x}$ and $y=k^x$ is $45^\circ$ is $\lambda$ then $|\ell n\lambda|$ is equal to
\qn{5} If the equation of tangent drawn to curve $y=f(x)$ at point P(3, 5) is $5x-4y+5=0$ and $\displaystyle\lim_{x\to3}\dfrac{\left(3^{4f(x)}-2\left(1+3+3^2+\cdots\cdots+3^{19}\right)-1\right)^2}{1-\cos(\ln(4-x))}=2\left(a\,3^{20}\,\ell n\,c\right)^2$ (where a and c are prime numbers), then a + c is
\qn{6} Let f : R $\to$ [$-$8, 8] be an onto function which is given by $f(x)=\dfrac{bx}{(a-3)x^3+x^2+4}$ \ a, b, $\in$ R$^+$. If the set of values of m for which the equation $f(x)=mx$ has 3 distinct real solution is (p, q), then the value of (a + b + p + q) is
\end{multicols}
\end{document}
